\documentclass[12pt]{article}
\usepackage{amsmath,amssymb,amsfonts}
\usepackage[margin=1in]{geometry}
\usepackage{bm}

\newcommand{\vect}[1]{\mathbf{#1}}
\newcommand{\grad}{\boldsymbol{\nabla}}

\begin{document}

\section*{Problem 7}

\textbf{Problem:} Prove Green's theorem, that is,
\[
\int_\tau (\varphi\,\nabla^2\psi - \psi\,\nabla^2\varphi)\,d\tau = \oint_\sigma (\varphi\,\grad\psi - \psi\,\grad\varphi)\cdot d\boldsymbol{\sigma}.
\]

\textit{Hint:} Apply Gauss's divergence theorem to the vectors $\vect{A} = \varphi\,\grad\psi$ and $\vect{B} = \psi\,\grad\varphi$.

\bigskip
\textbf{Solution:}

\subsection*{Step 1: Apply the divergence theorem to $\vect{A} = \varphi\,\grad\psi$}

Gauss's divergence theorem (Eq.~1.57 in Byron \& Fuller) states:
\[
\int_\tau \grad\cdot\vect{A}\;d\tau = \oint_\sigma \vect{A}\cdot d\boldsymbol{\sigma}.
\]

Compute $\grad\cdot(\varphi\,\grad\psi)$. Using the product rule for divergence:
\[
\grad\cdot(\varphi\,\grad\psi) = (\grad\varphi)\cdot(\grad\psi) + \varphi\,\nabla^2\psi.
\]

Applying the divergence theorem:
\begin{equation}
\int_\tau \bigl[(\grad\varphi)\cdot(\grad\psi) + \varphi\,\nabla^2\psi\bigr]\,d\tau = \oint_\sigma \varphi\,\grad\psi \cdot d\boldsymbol{\sigma}. \label{eq:green1}
\end{equation}

This is known as \textbf{Green's first identity}.

\subsection*{Step 2: Apply the divergence theorem to $\vect{B} = \psi\,\grad\varphi$}

By the same calculation with $\varphi$ and $\psi$ interchanged:
\begin{equation}
\int_\tau \bigl[(\grad\psi)\cdot(\grad\varphi) + \psi\,\nabla^2\varphi\bigr]\,d\tau = \oint_\sigma \psi\,\grad\varphi \cdot d\boldsymbol{\sigma}. \label{eq:green1b}
\end{equation}

\subsection*{Step 3: Subtract}

Subtract Eq.~\eqref{eq:green1b} from Eq.~\eqref{eq:green1}. The $(\grad\varphi)\cdot(\grad\psi)$ terms cancel (since the dot product is commutative):
\[
\int_\tau \bigl[\varphi\,\nabla^2\psi - \psi\,\nabla^2\varphi\bigr]\,d\tau = \oint_\sigma \bigl[\varphi\,\grad\psi - \psi\,\grad\varphi\bigr]\cdot d\boldsymbol{\sigma}.
\]

This is \textbf{Green's theorem} (also called Green's second identity). \hfill $\blacksquare$

\end{document}
